\documentclass{article}

% packages
\usepackage{amsmath, amsthm, thmtools, amsfonts, amssymb, luacode, catchfile, tikzducks, hyperref, ifthen}
\ifcsname c@kobocompile\endcsname
	\usepackage[a5paper, total={1072pt, 1448pt}, margin=10pt, includeheadfoot]{geometry} % set page margins
\else
	\usepackage[a4paper, margin=50pt, includeheadfoot]{geometry}
\fi
\usepackage[shortlabels]{enumitem}
\usepackage[skip=3pt, indent=0pt]{parskip}

% language
\usepackage[bidi=basic, layout=tabular, provide=*]{babel}
\ifcsname c@english\endcsname
	\babelprovide[main, import]{english}
\else
	\babelprovide[main, import]{hebrew}
	\babelprovide{rl}
\fi
%\babelfont{rm}{Libertinus Serif}
\babelfont{rm}[Renderer=Harfbuzz]{Libertinus Serif}
\babelfont{sf}{Libertinus Sans}
\babelfont{tt}{Libertinus Mono}

% style
\AddToHook{cmd/section/before}{\clearpage}	% Add line break before section
\linespread{1.3}
\setcounter{secnumdepth}{0}		% Remove default number tags from sections, this won't do well with theorems
\AtBeginDocument{\setlength{\belowdisplayskip}{3pt}}
\AtBeginDocument{\setlength{\abovedisplayskip}{3pt}}
\graphicspath{ {../images/} }

% operators
\DeclareMathOperator\cis{cis}
\DeclareMathOperator\Sp{Sp}
\DeclareMathOperator\tr{tr}
\DeclareMathOperator\im{Im}
\DeclareMathOperator\re{Re}
\DeclareMathOperator\diag{diag}
\DeclareMathOperator*\lowlim{\underline{lim}}
\DeclareMathOperator*\uplim{\overline{lim}}
\DeclareMathOperator\rng{rng}
\DeclareMathOperator\Sym{Sym}
\DeclareMathOperator\Arg{Arg}
\DeclareMathOperator\Log{Log}
\DeclareMathOperator\dom{dom}
\DeclareMathOperator\supp{Supp}
\DeclareMathOperator\var{Var}
\DeclareMathOperator\cov{Cov}

% commands
%\renewcommand\qedsymbol{\textbf{מש''ל}}
%\renewcommand\qedsymbol{\fbox{\emoji{lizard}}}
\newcommand{\Aa}[0]{\mathcal{A}}
\newcommand{\Bb}[0]{\mathcal{B}}
\newcommand{\CC}[0]{\mathbb{C}}
\newcommand{\Cc}[0]{\mathcal{C}}
\newcommand{\EE}[0]{\mathbb{E}}
\newcommand{\FF}[0]{\mathbb{F}}
\newcommand{\Ff}[0]{\mathcal{F}}
\newcommand{\Ii}[0]{\mathcal{I}}
\newcommand{\Gg}[0]{\mathcal{G}}
\newcommand{\Ll}[0]{\mathcal{L}}
\newcommand{\Mm}[0]{\mathcal{M}}
\newcommand{\NN}[0]{\mathbb{N}}
\newcommand{\Nn}[0]{\mathcal{N}}
\newcommand{\PP}[0]{\mathbb{P}}
\newcommand{\Pp}[0]{\mathcal{P}}
\newcommand{\QQ}[0]{\mathbb{Q}}
\newcommand{\RR}[0]{\mathbb{R}}
\newcommand{\Rr}[0]{\mathcal{R}}
\newcommand{\Ss}[0]{\mathcal{S}}
\newcommand{\TT}[0]{\mathbb{T}}
\newcommand{\Uu}[0]{\mathcal{U}}
\newcommand{\Vv}[0]{\mathcal{V}}
\newcommand{\Ww}[0]{\mathcal{W}}
\newcommand{\ZZ}[0]{\mathbb{Z}}
\newcommand{\acts}[0]{\circlearrowright}
\newcommand{\explain}[2] {
	\begin{flalign*}
		 && \text{#2} && \text{#1}
	\end{flalign*}
}
\newcommand{\maketitleprint}[0]{ \begin{center}
	%\begin{tikzpicture}[scale=3]
	%	\duck[graduate=gray!20!black, tassel=red!70!black]
	%\end{tikzpicture}	
	\includegraphics[width=6cm]{cover}
\end{center}
}

% theorem commands
\newtheoremstyle{c_remark}
	{}	% Space above
	{}	% Space below
	{}% Body font
	{}	% Indent amount
	{\bfseries}	% Theorem head font
	{}	% Punctuation after theorem head
	{.5em}	% Space after theorem head
	{\thmname{#1}\thmnumber{ #2}\thmnote{ \normalfont{\text{(#3)}}}}	% head content
\newtheoremstyle{c_definition}
	{3pt}	% Space above
	{3pt}	% Space below
	{}% Body font
	{}	% Indent amount
	{\bfseries}	% Theorem head font
	{}	% Punctuation after theorem head
	{.5em}	% Space after theorem head
	{\thmname{#1}\thmnumber{ #2}\thmnote{ \normalfont{\text{(#3)}}}}	% head content
\newtheoremstyle{c_plain}
	{3pt}	% Space above
	{3pt}	% Space below
	{\itshape}% Body font
	{}	% Indent amount
	{\bfseries}	% Theorem head font
	{}	% Punctuation after theorem head
	{.5em}	% Space after theorem head
	{\thmname{#1}\thmnumber{ #2}\thmnote{ \text{(#3)}}}	% head content

\ifcsname c@english\endcsname
	\theoremstyle{plain}
	\newtheorem{theorem}{Theorem}[section]
	\newtheorem{lemma}[theorem]{Lemma}
	\newtheorem{proposition}[theorem]{Proposition}
	\newtheorem*{proposition*}{Proposition}
	%\newtheorem{corollary}[theorem]{אין חלופה עברית}

	\theoremstyle{definition}
	\newtheorem{definition}[theorem]{Definition}
	\newtheorem*{definition*}{Definition}
	\newtheorem{example}{Example}[section]
	\newtheorem{exercise}{Exercise}[section]

	\theoremstyle{remark}
	\newtheorem*{remark}{Remark}
	\newtheorem*{solution}{Solution}
	\newtheorem{conclusion}[theorem]{Conclusion}
	\newtheorem{notation}[theorem]{Notation}
\else
	\theoremstyle{c_plain}
	\newtheorem{theorem}{משפט}[section]
	\newtheorem{lemma}[theorem]{למה}
	\newtheorem{proposition}[theorem]{טענה}
	\newtheorem*{proposition*}{טענה}
	%\newtheorem{corollary}[theorem]{אין חלופה עברית}

	\theoremstyle{c_definition}
	\newtheorem{definition}[theorem]{הגדרה}
	\newtheorem*{definition*}{הגדרה}
	\newtheorem{example}{דוגמה}[section]
	\newtheorem{exercise}{תרגיל}[section]

	\theoremstyle{c_remark}
	\newtheorem*{remark}{הערה}
	\newtheorem*{solution}{פתרון}
	\newtheorem{conclusion}[theorem]{מסקנה}
	\newtheorem{notation}[theorem]{סימון}
\fi

% Questions related commands
\newcounter{question}
\setcounter{question}{1}
\newcounter{sub_question}
\setcounter{sub_question}{1}

\ifcsname c@english\endcsname
	\newcommand{\question}[1][0]{
		\ifthenelse{#1 = 0}{}{\setcounter{question}{#1}}
		\section{Question \arabic{question}}
		\addtocounter{question}{1}
		\setcounter{sub_question}{1}
	}

	\newcommand{\subquestion}[1][0]{
		\ifthenelse{#1 = 0}{}{\setcounter{sub_question}{#1}}
		\subsection{Part \alph{sub_question}}
		\addtocounter{sub_question}{1}
	}
\else
	\newcommand{\question}[1][0]{
		\ifthenelse{#1 = 0}{}{\setcounter{question}{#1}}
		\section{שאלה \arabic{question}}
		\addtocounter{question}{1}
		\setcounter{sub_question}{1}
	}

	\newcommand{\subquestion}[1][0]{
		\ifthenelse{#1 = 0}{}{\setcounter{sub_question}{#1}}
		\subsection{סעיף \localecounter{letters.gershayim}{sub_question}}
		\addtocounter{sub_question}{1}
	}
\fi

% import lua and start of document
\directlua{common = require ('../common')}

\GetEnv{AUTHOR}

% headers
\author{\AUTHOR}
\date\today

\title{פתרון ממ''ן 15 – אלגברה לינארית 2 (20229)}

\begin{document}
\maketitle
\maketitleprint{}

\section{שאלה 1}
\subsection{סעיף א'}
תהי העתקה לינארית $T : V \to V$ אשר בבסיס הסטנדרטי מיוצגת על־ידי המטריצה
\[
	A = \begin{pmatrix}
		1 & 5 \\
		-10 & -1
	\end{pmatrix}
\]
נמצא את כל תת־המרחבים ה־$T$ שמורים של $V$ כאשר $V = \RR^2, \CC^2$: \\*
\textbf{(1)}
נגדיר $V = \RR^2$: \\*
תחילה, נמצא את ערכיה העצמיים של $T$ על־ידי חישוב פולינום אופייני:
\[
	p(t) = \begin{vmatrix}
		t - 1 & -5 \\
		10 & t + 1
	\end{vmatrix}
	= (t - 1)(t + 1) + 50
	= t^2 + 49
\]
אין להעתקה $T$ אם־כן ערכים עצמיים כלל, ולכן משאלה 8.4.3 א' נובע כי אין ל־$T$ תת־מרחב $T$ שמור שאיננו טריוויאלי. \\*
בהתאם כלל התת־מרחבים ה־$T$ שמורים הם, על־פי דוגמה 8.4.2, הם מרחב האפס ו־$V$ עצמו. \\*
\textbf{(2)}
נגדיר $V = \CC^2$: \\*
מעל שדה המרוכבים $p(t) = t^2 + 49 = (t - 7i)(t + 7i)$ ובהתאם $-7i, 7i$ הם ערכיה העצמיים של $T$ והיא אף לכסינה. \\*
נמצא וקטורים עצמיים של $T$ על־ידי חישוב המרחב העצמי לכל ערך עצמי:
\begin{align*}
	& (7i I - A) {(x, y)}^t = 0 \rightarrow
	\begin{pmatrix}
		-1 + 7i & -5 \\
		10 & 1 + 7i
	\end{pmatrix}
	\rightarrow
	\begin{pmatrix}
		-50 & -5(1 + 7i) \\
		10 & 1 + 7i
	\end{pmatrix}
	\rightarrow
	10x - (1 + 7i)y = 0
	\rightarrow
	\Sp\{ (1 + 7i, 10) \} \\
	& (-7i - A) {(x, y)}^t = 0
	\rightarrow
	\begin{pmatrix}
		-1 -7i & -5 \\
		10 & 1 - 7i
	\end{pmatrix}
	\rightarrow
	\begin{pmatrix}
		-50 & -5(1 - 7i) \\
		10 & 1 - 7i
	\end{pmatrix}
	\rightarrow
	10x + (1 - 7i)y = 0
	\rightarrow
	\Sp\{ (7i - 1, 10) \}
\end{align*}
אז מצאנו כי כל תת־המרחבים ה־$T$ שמורים שאינם טריוויאליים הם
\[
	\Sp\{ (7i - 1, 10) \},
	\Sp\{ (7i + 1, 10) \}
\]

\subsection{סעיף ב'}
יהי $V$ מרחב לינארי מעל $F$ ותהי $T : V \to V$ העתקה לינארית כך שכל תת־מרחב של $V$ הוא $T$ שמור. \\*
נוכיח שקיים $\alpha \in F$ כך ש־$T = \alpha I$.
\begin{proof}
	ידוע כי כל תת־מרחב של $V$ הוא $T$ שמור ולכן נובע כי גם עבור תת־מרחבים מממד $1$ הם $T$ שמורים. \\*
	מהגדרת תת־מרחב $T$ שמור נובע ישירות כי עבור צמצום למרחב מממד $1$ מקיים הצמצום $T_{\dim 1} u = \lambda u$ כאשר $\lambda \in F$. \\*
	אנו יכולים להסיק מסיבה זו שההעתקה $T$ היא לכסינה ונגדיר $\lambda_1, \hdots, \lambda_n$ ערכיה העצמיים. \\*
	יהיו $i, j$ כאשר $1 \le i < j \le n$. נגדיר $V_{ij} = \Sp\{ u_i + u_j \}$ תת־מרחב המוגדר על־ידי הווקטורים העצמיים המתאימים ל־$\lambda_i, \lambda_j$. \\*
	כמובן שתת־מרחב זה הוא $T$ שמור, ולכן $T(u_i + u_j) = T u_i + T u_j = \lambda_i u_i + \lambda_j u_j \in \Sp\{ u_i + u_j \}$.
	נובע כי $\lambda_i u_i + \lambda_j u_j = \alpha u_i + \alpha u_j$ ולכן $\lambda_i = \lambda_j = \alpha$.
	מצאנו כי $\lambda_i = \alpha$ לכל $1 \le i \le n$ ובהתאם $T u = \alpha u$ לכל $u \in V$.
\end{proof}

\section{שאלה 2}
\subsection{סעיף א'}
תהי $T$ העתקה לינארית במרחב לינארי $V$ שממדו סופי ויהי $W$ תת־מרחב $T$ שמור של $V$ ו־$T_W$ הצמצום של $T$ ל־$W$.
\textbf{(1)}
נוכיח כי הפולינום המינימלי של $T_W$ מחלק את הפולינום המינימלי של $T$.
\begin{proof}
	נגדיר $M_1(x)$ הפולינום המינימלי של $T$ ו־$M_2(x)$ הפולינום המינימלי של $T_W$. \\*
	ידוע כי $M_1(T) = M_2(T_W) = 0$. נניח בשלילה כי $M_1(T_W) \ne 0$, ולכן קיים $u \in W$ כך ש־$M_1(T_W) u \ne 0$.
	אבל $W \subseteq V$ ולכן $u \in V$ ובהתאם $M_1(T) u = 0$, וידוע כי מעל $W$ מתקיים $T = T_W$ ולכן $M_1(T_w) = 0$ בסתירה לטענה. \\*
	משאלה 9.9.1 סעיף א' נובע ישירות כי $M_2$ מחלק את $M_1$.
\end{proof}
\textbf{(2)}
נוכיח כי אם $T$ לכסינה אז גם $T_W$ לכסינה.
\begin{proof}
	ממשפט 10.2.11 נובע כי $M_1$ מתפרקת למכפלת גורמים לינאריים שונים וידוע לנו כי $M_2$ מחלקת אותה,
	לכן גם $M_2$ מורכבת ממכפלת גורמים לינאריים שונים. ממשפט 10.2.11 נובע מסיבה זו שגם $M_2$ לכסינה.
\end{proof}

\subsection{סעיף ב'}
נניח כי $T : \RR^3 \to \RR^3$ היא בעלת ערכים עצמיים $1, 2, 3$ ווקטורים עצמיים $v_1, v_2, v_3$ בהתאמה. \\*
נמצא את כל תת־המרחבים ה־$T$ שמורים של $\RR^3$: \\*
ההעתקה $T$ היא בעלת $3$ ערכים עצמיים שונים ולכן לכסינה. לכל וקטור עצמי $u$ שנבחר קיים ערך עצמי $\lambda$ המקיים $T u = \lambda u$,
ולכן הצמצום של $T$ ל־$\Sp\{ u \}$ הוא $T$ שמור.
כמובן שגם חיבור שני תתי־מרחב $T$ שמורים יובילו לתת־מרחב $T$ שמור, ולכן כלל התת־מרחבים ה־$T$ שמורים על $\RR^3$ הם:
\[
	0,
	\Sp\{ v_1 \},
	\Sp\{ v_2 \},
	\Sp\{ v_3 \},
	\Sp\{ v_1, v_2 \},
	\Sp\{ v_1, v_3 \},
	\Sp\{ v_2, v_3 \},
	\RR^3
\]

\section{שאלה 3}
תהי $T : \RR^3 \to \RR^3$ העתקה לינארית אשר מיוצגת בבסיס הסטנדרטי על־ידי המטריצה
\[
	A = \begin{pmatrix}
		3 & 1 & 0 \\
		0 & 3 & 0 \\
		0 & 0 & 2
	\end{pmatrix}
\]
\subsection{סעיף א'}
נמצא תתי־מרחב $T$ שמורים לא טריוויאליים על $\RR^3$. \\*
$A$ היא מטריצה משולשית ולכן $T$ ניתן לשילוש ומטענה 8.1.1 נובע כי ערכיה הפנימיים של $T$ הם $2, 3$.
מחישוב המרחב העצמי אנו למדים כי $T(0, 0, 1) = 2(0, 0, 1)$ וגם כי $T(1, 0, 0) = 3(1, 0, 0)$ ולכן משאלה 8.4.1 סעיף ג' נובע כי המרחבים
\[
	\Sp\{ (1, 0, 0) \},
	\Sp\{ (0, 0, 1) \}
\]
הם מרחבים לא טריוויאליים $T$ שמורים מעל $\RR^3$.

\subsection{סעיף ב'}
יהי $W = \ker(T - 3I)$, נוכיח כי לא קיים תת־מרחב $U \subseteq \RR^3$ שהוא $T$ שמור ומקיים
\[
	\RR^3 = W \oplus U
\]
\begin{proof}
	מחישוב מערכת המשוואות אנו מקבלים כי
	\[
		W = \Sp\{ (1, 0, 0) \}
	\]
	לכן ממשפט הממד והגדרת החיבור הישר נובע כי $\dim U = 2$ בלבד. \\*
	על־פי סעיף א' וטענה 8.4.16 תת־המרחב היחיד ה־$T$ שמור מממד $2$ הוא
	\[
		U = \Sp\{ (1, 0, 0), (0, 0, 1) \}
	\]
	אבל במצב זה $W \cap U \ne \{ 0 \}$ ולכן לא קיים תת־מרחב $U$ המקיים את הטענה.
\end{proof}

\section{שאלה 4}
יהי $V$ מרחב לינארי מממד סופי ו־$T : V \to V$ העתקה לינארית. \\*
יהי $M(x) = M_1(x) M_2(x) \cdots M_k(x)$ הפולינום המינימלי של $T$ כאשר $M_i(x)$ פולינומים מתוקנים זרים בזוגות. \\*
נגדיר $V = W_1 \oplus \cdots \oplus W_k$ הפירוק הפרימרי המתאים ל־$T$ כאשר $W_i = \ker M_i(T)$. \\*
יהי $W$ תת־מרחב $T$ שמור של $V$. \\*
נוכיח כי
\[
	W = (W \cap W_1) \oplus \cdots \oplus (W \cap W_k)
\]
\begin{proof}
	יהי $1 \le i, j \le k$, על־פי הפירוק הפרימרי $W_i \cap W_j = \{0\}$ ולכן
	\[
		W_i \cap W_j \cap W = (W_i \cap W) \cap (W_j \cap W) = \{ 0 \} \tag{1}
	\]
	כמו־כן
	\begin{align*}
		(W_i \cap W) + \left(W_j \cap W\right)
		& = \{ u_i + u_j \mid u_i \in (W_i \cap W), u_j \in (W_j \cap W) \} \\
		& = \{ u_i + u_j \mid u_i \in W_i \cap W, u_j \in W_j \cap W \} \cap W & \text{ידוע כי $W$ מרחב לינארי} \\
		& = \{ u_i + u_j \mid u_i \in W_i, u_j \in W_j \} \cap W & \text{כנביעה מתכונות המרחב} \\
		& = (W_i + W_j) \cap W
	\end{align*}
	אז בתהליך דומה נוכל לקבוע גם כי
	\[
		(W \cap W_1) + \cdots + (W \cap W_k) = W \cap (W_1 + \cdots + W_k) = W \cap V = W
	\]
	וכנביעה מ־$(1)$ גם
	\[
		W = (W \cap W_1) \oplus \cdots \oplus (W \cap W_k)
	\]
\end{proof}

\section{שאלה 5}
יהי $V$ מרחב אוניטרי מממד סופי ו־$T : V \to V$ העתקה נורמלית. \\*
נוכיח שכל תת־מרחב $T$ שמור הוא גם $T^*$ שמור.
\begin{proof}
	יהי $W$ תת־מרחב $T$ שמור של $V$ ונגדיר $T_W$ הצמצום של $T$ על $W$. \\*
	אנו יודעים כי עבור וקטורים ב־$W$ ההעתקות $T$ ו־$T_W$ זהות ולכן גם $T_W$ נורמלית. \\*
	בשל כך ועל־פי משפט 3.2.1 ההעתקה $T_W$ לכסינה אוניטרית. \\*
	ממשפט 3.4.2 נובע כי אם $\lambda_i$ כאשר $1 \le i \le k$ ערכים עצמיים של $T_W$ ו־$P_i$ ההיטלים האורתוגונליים של $V_{\lambda_i}$, אז לכל $w \in W$
	\[
		T_W w = \lambda_1 P_1 w + \cdots + \lambda_k P_k w \in W
	\]
	ובהתאם גם לכל $i$
	\[
		\lambda_i P_i w \in W
	\]
	מלמה 3.2.5 וממשפט 3.4.2 סעיף ה' נובע כי
	\[
		\overline{\lambda_i} P_i \in W
	\]
	ולכן
	\[
		T^* w = \overline{\lambda_1} P_1 w + \cdots + \overline{\lambda_k} P_k w \in W
	\]
	ומכאן נובע כי $W$ הוא תת־מרחב $T^*$ שמור.
\end{proof}

\end{document}
